\chapter{Implementation Details}\label{chapter:chap6}

\section{Technology Stack}\label{sect:tech-stack}

\begin{tabularx}{\textwidth}{|l|l|X|}
\hline
\textbf{Component} & \textbf{Technology} & \textbf{Justification} \\
\hline
Core application & Rust & Memory safety, performance, async and concurrent support \\
\hline
CLI framework & Rust + Clap & Easy command definition with subcommands \\
\hline
Async runtime & Tokio & Industry-standard Rust async runtime, speed and ease of use \\
\hline
Database ORM & Diesel (async) & Type-safe SQL queries, fastest ORM, migration management \\
\hline
Connection pooling & bb8 & Fastest async connection pool for PostgreSQL \\
\hline
HTTP client & reqwest & Easy async HTTP for LLM API communication \\
\hline
JSON parsing & sonic\_rs & High-performance JSON parser (faster than serde\_json) \\
\hline
Work queue & Redis & Minimal latency, atomic operations, Lua scripting \\
\hline
Data storage & PostgreSQL & Fast, Relational integrity, complex query support, JSONB \\
\hline
Web framework & Axum & Async, Ease of use, performance, Tower middleware ecosystem \\
\hline
Parallelism & Rayon & Easy parallelism for typosquatting analysis \\
\hline
Logging & tracing + tracing-subscriber & Structured logging with file and console output, helpful for debugging and information \\
\hline
Containerization & Docker + Docker Compose & Reproducible and easy deployment, horizontal scaling \\
\hline
\end{tabularx}

\section{CLI Application Structure}\label{sect:cli-structure}

The CLI app has the following commands and structure:

\begin{tabularx}{\textwidth}{|l|X|}
\hline
\textbf{Command} & \textbf{Description} \\
\hline
\texttt{update-database} & Downloads and imports the crates.io CSV dump into PostgreSQL \\
\hline
\texttt{run-analysis} & Constructs dependency graph and populates Redis queue \\
\hline
\texttt{worker} & Starts a worker process that consumes from the queue, used from Docker \\
\hline
\texttt{compute-dependencies} & Post-processing for dependency risk flags \\
\hline
\texttt{queue-status} & Displays Redis queue statistics \\
\hline
\texttt{recover-stale} & Recovers items from crashed workers \\
\hline
\texttt{clear-queue} & Resets the entire queue state \\
\hline
\texttt{fix-cycles} & Detects and resolves dependency graph cycles \\
\hline
\texttt{scan-crate --name X} & Single-crate analysis with full dependency tree \\
\hline
\texttt{queue-crate --name X} & Queues a single crate and its transitive dependencies \\
\hline
\texttt{check-llm} & Verifies LM Studio connectivity \\
\hline
\texttt{typosquat} & Runs typosquatting detection on the entire database \\
\hline
\end{tabularx}

\section{Redis Queue Implementation}\label{sect:redis-queue}

The Redis queue implementation was designed to handle multiple complex scenarios:

- Topological ordering of crates to ensure the dependents are analyzed first

- Atomic operations - Atomic claim of crates, handling crate completion, dependent updates and others happen in an atomic manner to prevent race conditions

- Blocking pop - BRPOP is used to avoid busy waiting

- Statistics tracking - The queue has counters for total crates, processed crates and others to provide monitoring capabilities.

\section{Cycle Detection and Resolution}\label{sect:cycle-detection}

Dependency cycles can prevent crates from ever being processed as dependency cycles will cause crates to always have dependencies.

- Cycle detection - Using depth-first search, we identify cycle edges

- Cycle breaking - Removal of edges allows crates to be processed, 'fix-cycles' is used for simple A-B-A cycles, and 'break-cycles' is used for more complex cycles.

\section{Worker Lifecycle and Fault Tolerance}\label{sect:worker-lifecycle}

Each worker is an independent instance with a lifecycle consists of;

- Initialization - Connects to PostgreSQL and Redis and sets up the state and handlers

- Main loop - Claim a crate from the ready queue in an atomic way, download its repository and run each module, after that insert the results into PostgreSQL and remove the crate from Redis.

- Shutdown - Handle signals like SIGINT and SIGTERM by finishing the current crate and exit only after that.

\section{Batch Processing and Database Optimization}\label{sect:batch-db-opt}

- Batch updating - INSERT and on conflict UPDATE statements for results

- Chunked imports - The database population from the CSV files is done in chunks to improve performance and reduce memory usage

- Connection pooling - bb8 is used to manage database connections

- Indexing - Indexes on commonly used queries in the codebase are used; they do not increase the database size enough to be a problem